\documentclass{beamer}
\usepackage{tikz}
\usetheme{Berkeley}
\usecolortheme{default}

\title{Résolution algorithmique du jeu \url{loopover.xyz}} % il y a une option je sais pas pourquoi
% \subtitle{...}
\author{Parent Quentin}

\institute[VFU]{} % on verra
\date[VLC 2021]{} % à ajouter

\logo{%??? on verra
}
\AtBeginSection[]
{
	\begin{frame}{Table des matières}
		\tableofcontents[currentsection]
	\end{frame}
}

\begin{document}
	\frame{\titlepage}
	
	\begin{frame}{Jeu}
		insérer ici un joli dessin (deux, après un U1?)		
	\end{frame}
	
	\section{Approche calculatoire}
	\begin{frame}{Algorithme}
		description avec des jolis dessins
	\end{frame}
	
	\begin{frame}{Majoration}
		faire le calcul en fonction de $N$ puis conclure à $125$? (vérif)
	\end{frame}
	
	\begin{frame}{Minoration}
		Nombre de mélanges en moins de $k$ mouvements, pour un puzzle de taille $N$:
		$$ 1+\sum_{i=0}^k (4N-1)^i = 1 + \frac{1-(4N-1)^{k+1}}{2-4N}$$
		Si $k$ est la majoration recherchée, on a donc:
		\begin{align*}
			& \frac{N^2!}2  \geqslant 1 + \frac{1-(4N-1)^{k+1}}{2-4N} \\
			\iff & k \geqslant \frac{\log\left(\left(\frac{N^2!}2-1\right)(4N-2)-1\right)}{\log(4N-1)}
		\end{align*}
		ie, pour $N=5$, $k\geqslant 20$.
	\end{frame}
	
	
	\section{Résolution}
	\begin{frame}{Approche naïve}
		bfs!!! c'est long...
	\end{frame}
	\begin{frame}{IDA*}
		ben c'est un IDA* quoi
	\end{frame}
	\begin{frame}{Résolution par phases}
		préciser que réduit au groupe engendré par blabla $\rightarrow$ nice
		
		théorie des groupes: sous-groupes normaux 
	\end{frame}
	
	\begin{frame}{Comparaison}
		\begin{center}
			faire un joli tableau temps/nb mvts sur bruteforce/IDA*/deux phases/trois phases
%			\begin{tabular}{||c c c >+< c||} 
%				\hline
%				Col1 & Col2 & Col2 & Col3 \\ [0.5ex] 
%				\hline
%				1 & 6 & 87837 & 787 \\ 
%				\hline
%				2 & 7 & 78 & 5415 \\
%				\hline
%				3 & 545 & 778 & 7507 \\
%				\hline
%				4 & 545 & 18744 & 7560 \\
%				\hline
%				5 & 88 & 788 & 6344 \\ [1ex] 
%				\hline
%			\end{tabular}
		\end{center}
	\end{frame}
	

	\section{Exploration}
	\begin{frame}{Exploration exhaustive} % voir si je le mets en annexe?
		quel goat ce Heap
		
		(pour injection de $[\![1, k]\!]$ dans $[\![1, n]\!]$, Heap+$\binom n k$)
	\end{frame}
	\begin{frame}{Réduction par symétries}
		je sais pas ce que je dis ici, on prie??
	\end{frame}

	\begin{frame}{Résultats}
		on fait un joli récap
	\end{frame}
		
	\section{Conclusion}
	\begin{frame}{Résolution monophase}
		Les mélanges longs		
		
		et ça nous donne une minoration :D
	\end{frame}
	
	\begin{frame}{Bibliographie}
		ben euh la biblio
	\end{frame}
	
	\section{Annexes}
	\begin{frame}{euhhhh}
		on verra
	\end{frame}
\end{document}